% Section fragment; input from the main document.
\section{问题重述}
\label{sec:frontmatter-restatement}

\subsection{问题背景}

微网将分布式光伏、居民用电与储能连接起来，为清洁能源的本地消纳提供了途径。然而，光伏出力与用电需求存在时间错位，外网电价也有峰谷差异。因此，非孤岛式微网需要联合安排外网购电与储能充放电，在保障供电的同时降低购电费用，发挥储能的经济价值。

\subsection{问题要求}

依据典型日数据、全年负载与光伏实测数据、分时光伏预报及全年电价，结合储能参数，需解决以下问题：

\textbf{问题一：}在给定光伏预测及每日电价、负载相同的条件下，于0:00制定全天购电与充放电方案，满足负载、设备限制及日初日末储电量相等的要求，并计算购电费用。

\textbf{问题二：}在电价固定、负载与光伏随时间变化的条件下，于每天0:00制定购电计划；供电不足时按交易电价的5倍紧急购电，其余费用按计划量结算，求取相应购电与储能策略。

\textbf{问题三：}利用每天0:00、6:00、12:00和18:00发布的未来24小时光伏预报，制定并调整购电计划；下调部分按交易电价的50\%支付违约费，上调部分按1.5倍计价，综合计划、调整与紧急购电费用，分析日内修正的必要性。

\textbf{问题四：}将固定电价扩展为实时波动电价，重新计算问题二、三，比较不同信息条件与调整方式下的购电、储能策略及运行费用。
